\documentclass{article}

\usepackage{amsmath,amssymb}
\usepackage{xurl}
\usepackage[hidelinks]{hyperref}

% Shared commands for notes in docs/paper-gaps/.

\DeclareUnicodeCharacter{2080}{\ifmmode{}_{0}\else\textsubscript{0}\fi}
\DeclareUnicodeCharacter{2081}{\ifmmode{}_{1}\else\textsubscript{1}\fi}
\DeclareUnicodeCharacter{2082}{\ifmmode{}_{2}\else\textsubscript{2}\fi}
\DeclareUnicodeCharacter{2099}{\ifmmode{}_{n}\else\textsubscript{n}\fi}

\newcommand{\Lean}{\textsc{Lean}}
\ExplSyntaxOn
\NewDocumentCommand{\leanid}{m}{
  \ifmmode
    \textsf{\detokenize{#1}}
  \else
    \group_begin:
    \urlstyle{sf}
    \tl_set:Nn \l_tmpa_tl {#1}
    \tl_replace_all:Nnn \l_tmpa_tl {\_} {_}
    \exp_args:NV \path \l_tmpa_tl
    \group_end:
  \fi
}
\ExplSyntaxOff
\providecommand{\tr}{\operatorname{tr}}
\newcommand{\tnleanIssueBase}{https://github.com/LionSR/TNLean/issues/}
\newcommand{\tnleanPrBase}{https://github.com/LionSR/TNLean/pull/}
\newcommand{\tnleanDocBase}{https://sirui-lu.com/TNLean/docs/}
\newcommand{\ghissue}[1]{\href{\tnleanIssueBase#1}{GitHub issue \##1}}
\newcommand{\ghpr}[1]{\href{\tnleanPrBase#1}{PR \##1}}
\newcommand{\doclink}[1]{\href{\tnleanDocBase#1.html}{\path{#1}}}

% Dirac notation and the identity operator, as in blueprint/src/macros/common.tex.
% \providecommand keeps note-local definitions authoritative.
\usepackage{dsfont}
\providecommand{\ket}[1]{|#1\rangle}
\providecommand{\bra}[1]{\langle #1|}
\providecommand{\braket}[2]{\langle #1 | #2 \rangle}
\providecommand{\ketbra}[2]{|#1\rangle\!\langle #2|}
\providecommand{\Id}{\mathds{1}}
\providecommand{\one}{\mathds{1}}
\providecommand{\C}{\mathbb{C}}
\providecommand{\MN}[1]{M_{#1}(\C)}
\providecommand{\MD}{M_D(\C)}

% External referencing for paper sources.  The build helper writes sanitized
% auxiliary files under build/paper-aux/ and excludes duplicated source labels.
\usepackage{xr}

% Import a generated paper auxiliary file with a caller-supplied label prefix.
% The candidate paths support builds from docs/paper-gaps/, the repository root,
% and build/paper-aux/ itself.
\newcommand{\importpaperaux}[2]{%
  \IfFileExists{../../build/paper-aux/#2.aux}{%
    \externaldocument[#1]{../../build/paper-aux/#2}%
  }{%
    \IfFileExists{build/paper-aux/#2.aux}{%
      \externaldocument[#1]{build/paper-aux/#2}%
    }{%
      \IfFileExists{#2.aux}{%
        \externaldocument[#1]{#2}%
      }{}%
    }%
  }%
}

% General external reference macros
\newcommand{\paperref}[2]{\ref{#1-#2}}
\newcommand{\papereqref}[2]{(\ref{#1-#2})}

% CPSV16 source (1606.00608 / MPDO-22-12-17-2.tex) external document and shortcuts
\importpaperaux{cpsv16-}{1606.00608/MPDO-22-12-17-2-tnlean}
\newcommand{\cpsvref}[1]{\paperref{cpsv16}{#1}}
\newcommand{\cpsveqref}[1]{\papereqref{cpsv16}{#1}}

% Verdict marker: \gapnote{<kind>}{<status>}. Kind is the policy
% classification (clarification, local-correction, scope-restriction,
% unfaithful, false-source, open-gap); status is open, wip (elimination
% underway), resolved, or historical. Machine-read by the site index;
% prints a verdict line.
\newcommand{\gapnote}[2]{\par\noindent\textbf{Verdict:} #1 (#2).\par}


\title{Specified-Tensor Shift Values and the Public MPU Index}
\author{}
\date{2026-08-30}

\begin{document}
\maketitle
\gapnote{scope-restriction}{open}

\section{Source statement}

For a simple matrix product unitary tensor with positive right and left
source-cut ranks $r$ and $\ell$, Definition~IV.1 of
Ref.~\cite{Cirac2017MPU} uses
\begin{align}
    \frac12(\log_2 r-\log_2\ell).
    \label{eq:mpu_shift_specified_tensor_index_scope_expression}
\end{align}
Proposition~IV.2 proves that this value is independent of the chosen simple
blocking. The resulting blocking-independent quantity is the public MPU
index.
% Source: paper_v2.tex lines 681--704, labels def:index and
% index-well-defined.

For the topological-insulator examples, the paper states that the right shift
$T^{(N)}$ has index $\log_2 d$, its adjoint left shift $T^{(N)\dagger}$ has
index $-\log_2 d$, and the three displayed families
\begin{align}
    U_1^{(N)} &= (\Id\otimes\Id)^{\otimes N}, \\
    U_2^{(N)} &= T^{(N)\dagger}\otimes T^{(N)}, \\
    U_3^{(N)} &= T^{(N)}\otimes T^{(N)\dagger}
\end{align}
have index zero.
% Source: paper_v2.tex lines 1980--2001 and 2037--2041, label threeMPU.

\section{Formal boundary}

The definition \leanid{MPOTensor.sourceIndexValue} evaluates
Eq.~\ref{eq:mpu_shift_specified_tensor_index_scope_expression} for one named
tensor and supplied positivity proofs. It does not choose a simple blocking or
prove independence of that choice.

The shift calculations use the exact source-rank pairs
\begin{align}
    (r[T],\ell[T])
        &=(d^2,1), \\
    (r[T^\dagger],\ell[T^\dagger])
        &=(1,d^2).
\end{align}
The bond-one identity has equal ranks $(d,d)$. Source ranks multiply under the
independent tensor product, so each specified tensor for $U_1,U_2,U_3$ has
equal right and left rank $d^2$. Therefore the formal declarations
\leanid{MPOTensor.sourceIndexValue_rightShiftTensor},
\leanid{MPOTensor.sourceIndexValue_leftShiftTensor}, and the four corresponding
zero-value theorems prove the exact formulas for these named tensors.

These calculations verify the numerical expression and its signs. They do not
yet identify those values with the public blocking-independent MPU index
asserted by the source.

For the right and left shifts the choice of blocking is no longer a
restriction. For every $k\ge0$, the right shift blocked over $k+1$ sites is
simple and has source ranks
\begin{align}
    (r,\ell)=(d^{k+2},d^{k}),
\end{align}
the left shift has the reversed pair, and the blocked tensors for $U_2$ and
$U_3$ have $r=\ell=d^{2k+2}$. The declarations
\leanid{MPOTensor.sourceIndexValue_blockTensor_rightShiftTensor},
\leanid{MPOTensor.sourceIndexValue_blockTensor_leftShiftTensor},
\leanid{MPOTensor.sourceIndexValue_blockTensor_shiftExampleU₂}, and
\leanid{MPOTensor.sourceIndexValue_blockTensor_shiftExampleU₃} give the values
$\log_2 d$, $-\log_2 d$, $0$, and $0$ at every such blocking. The left shift
blocked over $k+1$ sites is also simple, but simplicity of the blockings of
$U_2$ and $U_3$ is not proved, so for these two tensors the zero values are
values of the source-index expression, not yet the index of Definition IV.1.
What remains is the step from these displayed tensors to an arbitrary tensor
generating the same operators, which needs the canonical form of the shift and
its uniqueness.

The review states that the shift cannot be approximated by a short-range time
evolution (arXiv:2011.12127, lines 2607--2611, citing the source). Unequal
source ranks alone do not give this: after blocking, every MPU, the shift
included, is a depth-two circuit of nearest-neighbour unitaries in its standard
form (Definition SF, lines 605--622, and the appendix, line 2308). The
statement that the shift is not a finite-depth circuit of local gates needs
three further results, none of which is formalized: additivity of the index
under composition, Theorem IndexTh (ii), lines 824--847; index zero for a
single layer of nearest-neighbour gates; and the robustness and equivalence
parts (iii)--(iv) of the same theorem.

\section{Elimination plan}

The public index construction is tracked by \ghissue{7011}. It must choose a
simple blocking and prove that
Eq.~\ref{eq:mpu_shift_specified_tensor_index_scope_expression} is independent
of every such choice. Its remaining source-faithful input is the four-way
source theorem tracked by \ghissue{5856}, which supplies the rank-product
identity needed by the blocking argument without adding it as a hypothesis.

After that construction, the shift results can be stated for the public index
by evaluating it at the displayed simple representatives and applying
blocking independence. Until then, the Lean declarations and their Blueprint
entry remain explicitly restricted to specified tensors.

\section{Verdict}

The specified-tensor formulas agree exactly with the source-rank calculations,
including the positive sign for the right shift and the negative sign for the
left shift. The open gap is only the passage from these named-tensor values to
the public blocking-independent MPU index of Definition~IV.1 and
Proposition~IV.2 of Ref.~\cite{Cirac2017MPU}.

\bibliographystyle{alpha}
\bibliography{references}

\end{document}